Le filament de tungstène d'une ampoule a une résistance de \qty{\v{R0}}{\ohm} à \qty{20}{\celsius}. Le coefficient de température du tungstène est de \qty{4.5e-3}{\celsius^{-1}}. En fonctionnement, le filament atteint \qty{\v{T}}{\celsius} et l'ampoule est branchée sous \qty{120}{V}.

\begin{enumerate}[label=\alph*)]
    \item Déterminez la résistance du filament à sa température de fonctionnement (approximation linéaire).
    \item Déterminez le courant et la puissance de l'ampoule en fonctionnement.
    \item Déterminez le courant au moment de l'allumage, quand le filament est encore à \qty{20}{\celsius}. Pourquoi les ampoules à filament grillent-elles le plus souvent à l'allumage ?
\end{enumerate}
